\section{Trabalhos Relacionados}
\label{sec:related}

As pesquisas sobre comunicação segura em IoT investigam a redução do custo de estabelecimento de sessões, a proteção do estado de retomada e o reconhecimento de operações repetidas. Nesta seção, são apresentados os principais trabalhos relacionados a esses mecanismos e suas diferenças em relação à recuperação de sessões e leituras avaliada neste artigo.

Avoine, Canard e Ferreira (2019) \cite{avoine2019iotfriendly} apresentam uma construção com retomada e sigilo futuro em uma arquitetura de três participantes, utilizando operações simétricas entre dispositivo e servidor. O protocolo evolui segredos e prevê ressincronização após mensagens perdidas durante sua atualização. Esse tratamento constitui um antecedente para a recuperação da comunicação em IoT. Nosso estudo examina a restauração de estado do cliente e o custo da recuperação da telemetria com TLS nativo. A evolução de segredos desse trabalho possui condições de segurança diferentes das do PSK estático utilizado em um dos controles do protótipo.

Catoni, Puliafito e Dini (2025) \cite{catoni2025continuity} investigam a continuidade de serviços em ambientes de borda e nuvem, combinando TLS 1.3, X25519MLKEM768 e retomada de sessões. A implementação utiliza Envoy e BoringSSL, com uma chave de proteção dos tickets compartilhada entre proxies, permitindo retomada sem cache por sessão. O estudo compara estabelecimento completo e retomado com troca clássica ou híbrida, medindo custos criptográficos, bytes e tempo de handshake em cenários LAN e LTE. A autenticação pós-quântica é indicada como trabalho futuro. Nosso estudo compara a preservação e a perda da chave dos tickets após reinícios do gateway, variando também a frequência de confirmações perdidas. A avaliação inclui autenticação pós-quântica e acompanha o tráfego até a confirmação da leitura, incluindo novas tentativas e efeitos no banco.

Hebrok et al. (2023) \cite{hebrok2023tickets} examinam riscos criptográficos em tickets TLS, analisando doze implementações e realizando medições em servidores públicos. O trabalho identifica problemas no gerenciamento das chaves que protegem os tickets e na reutilização de fluxos de cifragem. Esses resultados evidenciam a importância da proteção do estado de retomada. Nossa avaliação considera sua disponibilidade após reinícios e quantifica o custo de apresentar um ticket recusado, seguido de autenticação completa, em diferentes configurações criptográficas.

Lee et al. (2015) \cite{lee2015rifl} apresentam RIFL, um mecanismo para reconhecer chamadas remotas repetidas e conservar seus resultados. O registro de conclusão e os efeitos da operação são salvos juntos, de modo que ambos sejam gravados ou nenhum deles seja. Os dados permanecem disponíveis após reinícios, permitindo responder a novas tentativas sem repetir a execução. Esse princípio fundamenta o banco de leituras já processadas deste trabalho. Nosso estudo observa os efeitos registrados no gateway quando a mesma leitura é reenviada por outra sessão autenticada, após a perda da confirmação, comparando a supressão de duplicações ativa e desativada.

A contribuição deste estudo concentra-se na avaliação conjunta das dependências de estado da recuperação e de seus custos. O protótipo instrumentado observa os efeitos da restauração de chaves e contadores. Os estudos de caso com TLS nativo comparam configurações clássica, híbrida e pós-quântica, mantendo a mesma construção de tickets nas políticas com chave preservada ou perdida. A variação da frequência de reinícios e de confirmações perdidas permite relacionar o custo por leitura confirmada às falhas e às informações disponíveis. O histórico da aplicação é tratado separadamente, distinguindo a recuperação da sessão da supressão de efeitos repetidos no banco.
